\subsection{可分解代数的分解}

若 g 是 \(\mathfrak{gl}(V)\) 的一个李代数（具有根基 r），则 r 的幂零元素所成的集是 g 的一个幂零理想，即 g 的恒等表示的最大幂零性理想（第一章，§5，第 3 节，定理 1 的推论 6）。在本节中，我们将把这个理想记作 \(\mathfrak{n}_{\mathrm{V}}(\mathfrak{g})\)。它包含 g 的幂零根基 \([g,g]\cap r\)（第一章，§5，第 3 节，定理 1）。

命题 5. 设 g 是 \(\mathfrak{gl}(V)\) 的一个可分解幂零李子代数。设 t 是 g 的半单元素所成的集。则 t 是 g 的一个中心子代数，且 g 作为李代数是 t 与 \(\mathfrak{n}_{\mathrm{V}}(\mathfrak{g})\) 的积。

若 \(x \in t\)，则 \(ad_{g}x\) 既半单又幂零，故为零，从而 x 在 g 中是中心的。因而 t 是 g 的一个理想，且 \(t \cap \mathfrak{n}_{\mathrm{V}}(\mathfrak{g}) = 0\)。由于 g 可分解，\(\mathfrak{g} = \mathfrak{t} + \mathfrak{n}_{\mathrm{V}}(\mathfrak{g})\)，命题得证。

命题 6. 设 g 是 \(\mathfrak{gl}(\mathrm{V})\) 的一个可分解李子代数。设 T 是由半单元素组成的 g 的交换子代数所成的集，\(T_{1}\) 是 T 的极大元素所成的集。设 H 是 g 的嘉当子代数所成的集。

\begin{enumerate}
\def\labelenumi{(\roman{enumi})}
\item
  对 \(\mathfrak{h} \in \mathcal{H}\)，设 \(\varphi(\mathfrak{h})\) 是 \(\mathfrak{h}\) 的半单元素所成的集。则 \(\varphi(\mathfrak{h}) \in \mathcal{T}_1\)。
\item
  对 \(\mathfrak{t} \in \mathcal{T}_1\)，设 \(\psi(\mathfrak{t})\) 是 \(\mathfrak{t}\) 在 \(\mathfrak{g}\) 中的交换子。则 \(\psi(\mathfrak{t}) \in \mathscr{H}\)。
\item
  映射 \(\varphi\) 与 \(\psi\) 分别是从 \(\mathcal{H}\) 到 \(\mathcal{T}_1\) 与从 \(\mathcal{T}_1\) 到 \(\mathcal{H}\) 的互逆双射。
\item
  若 \(k\) 代数闭，则 \(\mathrm{Aut}_e(\mathfrak{g})\) 传递地作用于 \(\mathcal{T}_1\)。
\end{enumerate}

设 \(\mathfrak{h} \in \mathcal{H}\)，并置 \(\mathfrak{t} = \varphi(\mathfrak{h})\)。由命题 5 与命题 3 的推论 2，\(\mathfrak{t} \in \mathcal{T}\) 且 \(\mathfrak{h} = \mathfrak{t} \times \mathfrak{n}_{\mathrm{V}}(\mathfrak{h})\)。对任何子代数 \(\mathfrak{u}\)（\(\mathfrak{g}\) 的子代数），用 \(\psi(\mathfrak{u})\) 表示 \(\mathfrak{u}\) 在 \(\mathfrak{g}\) 中的交换子。则 \(\mathfrak{h} \subset \psi(\mathfrak{t})\)，且 \(\psi(\mathfrak{t}) \subset \mathfrak{g}^0(\mathfrak{h})\)，因为 \(\mathfrak{n}_{\mathrm{V}}(\mathfrak{h})\) 的元素都是幂零的，故 \(\mathfrak{h} = \psi(\mathfrak{t})\)。若 \(\mathfrak{t}' \in \mathcal{T}\) 且 \(\mathfrak{t} \subset \mathfrak{t}'\)，则 \(\mathfrak{t}' \subset \psi(\mathfrak{t}) = \mathfrak{h}\)，于是 \(\mathfrak{t}' = \mathfrak{t}\)，因而 \(\mathfrak{t} \in \mathcal{T}_1\)。

设 \(\mathfrak{t} \in \mathcal{T}_1\)，并置 \(\mathfrak{c} = \psi(\mathfrak{t})\)。设 \(\mathfrak{h}\) 是 \(\mathfrak{c}\) 的一个嘉当子代数。由 §2 第 3 节命题 10，\(\mathfrak{h} \in \mathcal{H}\) 且 \(\mathfrak{t} \subset \mathfrak{h}\)。置 \(\mathfrak{t}_1 = \varphi(\mathfrak{h}) \in \mathcal{T}\)。则 \(\mathfrak{t} \subset \mathfrak{t}_1\)，故 \(\mathfrak{t} = \mathfrak{t}_1\)，且由前所述 \(\mathfrak{h} = \psi(\mathfrak{t}_1) = \psi(\mathfrak{t}) = \mathfrak{c}\)。于是 \(\psi(\mathfrak{t}) \in \mathcal{H}\)，且 \(\varphi(\psi(\mathfrak{t})) = \mathfrak{t}\)。

我们就这样证明了 (i)、(ii) 与 (iii)。设 \(k\) 代数闭。由于 \(\mathrm{Aut}_e(\mathfrak{g})\) 传递地作用于 \(\mathcal{H}(\S 3, \text{no. } 2, \text{Th. } 1)\)，故 \(\mathrm{Aut}_e(\mathfrak{g})\) 传递地作用于 \(\mathcal{T}_1\)。

推论 1. \(\mathfrak{g}\) 的嘉当子代数是 \(\mathfrak{g}\) 的正则半单元素的中心化子。

若 \(x \in \mathfrak{g}\) 正则，则 \(\mathfrak{g}^0(x)\) 是 \(\mathfrak{g}\) 的一个嘉当子代数（§2，第 3 节，定理 1 (i)）；此外，若 \(x\) 半单，则 \(\mathfrak{g}^0(x)\) 是 \(x\) 在 \(\mathfrak{g}\) 中的中心化子。反之，设 \(\mathfrak{h}\) 是 \(\mathfrak{g}\) 的一个嘉当子代数。存在 \(\mathfrak{t} \in \mathcal{T}_1\)，使得 \(\mathfrak{h} = \psi(\mathfrak{t})\)。由 §1 第 2 节命题 7，存在 \(x \in \mathfrak{t}\)，使得 \(\mathfrak{h} = \mathfrak{g}^0(x)\)；由于 \(x \in \mathfrak{t}\)，\(\mathfrak{g}^0(x) = \mathfrak{g}_0(x)\)。由 §3 第 3 节定理 2 (ii)，\(x\) 正则。

推论 2. 再设 g 可解。则：

\begin{enumerate}
\def\labelenumi{(\roman{enumi})}
\item
  \(\operatorname{Aut}(\mathfrak{g})\) 中由诸 \(e^{\mathrm{ad}x}, x \in \mathcal{C}^{\infty}\mathfrak{g}\)（参见 §3 第 4 节）组成的子群传递地作用于 \(\mathcal{T}_1\)。
\item
  若 \(\mathfrak{t} \in \mathcal{T}_1\)，则 \(\mathfrak{g}\) 是 \(\mathfrak{t}\) 与 \(\mathfrak{n}_{\mathrm{V}}(\mathfrak{g})\) 的半直积。
\end{enumerate}

断言 (i) 由下列事实得出：诸 \(e^{\mathrm{ad}x}\)（\(x \in \mathcal{C}^{\infty}\mathfrak{g}\)）所成的群传递地作用于 \(\mathcal{H}(\S 3\)，第 4 节，定理 3）。

我们证明 (ii)。设 \(t \in T_{1}\)，并设 \(\mathfrak{h} = \psi(t)\) 是 g 的相应的嘉当子代数。鉴于命题 5，\(\mathfrak{h} = \mathfrak{t} + \mathfrak{n}_{\mathrm{V}}(\mathfrak{h}) \subset \mathfrak{t} + \mathfrak{n}_{\mathrm{V}}(\mathfrak{g})\)。另一方面，\(g = h + [g, g]\)（§2，第 1 节，命题 4 的推论 3）且 \([\mathfrak{g}, \mathfrak{g}] \subset \mathfrak{n}_{\mathrm{V}}(\mathfrak{g})\)，故 \(\mathfrak{g} = \mathfrak{t} + \mathfrak{n}_{\mathrm{V}}(\mathfrak{g})\)。但显然 \(\mathfrak{t} \cap \mathfrak{n}_{\mathrm{V}}(\mathfrak{g}) = \{0\}\)。于是代数 g 是 t 与理想 \(\mathfrak{n}_{\mathrm{V}}(\mathfrak{g})\) 的半直积。

命题 7. 设 \(\mathfrak{g}\) 是 \(\mathfrak{gl}(\mathrm{V})\) 的一个可分解李子代数。

\begin{enumerate}
\def\labelenumi{(\roman{enumi})}
\item
  存在一个李子代数 \(\mathfrak{m}\)（\(\mathfrak{g}\) 的李子代数），它在 \(\mathfrak{gl}(\mathrm{V})\) 中约化，使得 \(\mathfrak{g}\) 是 \(\mathfrak{m}\) 与 \(\mathfrak{n}_{\mathrm{V}}(\mathfrak{g})\) 的半直积。
\item
  \(\mathfrak{g}\) 的任何两个具有 (i) 中性质的李子代数都在 \(\mathrm{Aut}_e(\mathfrak{g})\) 下共轭。
\end{enumerate}

g 的根基 r 可分解（第 2 节，命题 4 的推论 2）。由命题 6 的推论 2，存在 r 的一个由半单元素组成的交换子代数 t，使得 \(\mathfrak{r} = \mathfrak{t} \oplus \mathfrak{n}_{\mathrm{V}}(\mathfrak{r})\)。由于 \(ad_{g}t\) 由半单元素组成，g 是 \([t, g]\) 与 t 的中心化子 z 的直和（第一章，§3，第 5 节，命题 6）。由于 \([t, g] \subset r\)，\(g = z + r\)。因而，若 s 是 z 的一个 Levi 子代数（第一章，§6，第 8 节），则 \(g = s + r\)，故 s 是 g 的一个 Levi 子代数。置 \(m = s \oplus t\)。由于 \([s, t] = \{0\}\)，m 是 g 的一个李子代数，且由第一章 §6 第 5 节定理 4，它在 \(\mathfrak{gl}(V)\) 中约化。此外，

\[
\mathfrak {g} = \mathfrak {s} \oplus \mathfrak {r} = \mathfrak {s} \oplus \mathfrak {t} \oplus \mathfrak {n} _ {\mathrm{V}} (\mathfrak {r}) = \mathfrak {s} \oplus \mathfrak {t} \oplus \mathfrak {n} _ {\mathrm{V}} (\mathfrak {g}) = \mathfrak {m} \oplus \mathfrak {n} _ {\mathrm{V}} (\mathfrak {g})
\]

因为 \(\mathfrak{n}_{\mathrm{V}}(\mathfrak{g}) = \mathfrak{n}_{\mathrm{V}}(\mathfrak{r})\)。故得 (i)。

现在设 \(\mathfrak{m}'\) 是 \(\mathfrak{g}\) 的一个李子代数，它是 \(\mathfrak{n}_{\mathrm{V}}(\mathfrak{g})\) 的补，且在 \(\mathfrak{gl}(\mathrm{V})\) 中约化。我们证明 \(\mathfrak{m}'\) 与 \(\mathfrak{m}\) 在 \(\operatorname{Aut}_e(\mathfrak{g})\) 下共轭。我们有 \(\mathfrak{m}' = \mathfrak{s}' \oplus \mathfrak{t}'\)，其中 \(\mathfrak{s}' = [\mathfrak{m}', \mathfrak{m}']\) 半单，且中心 \(\mathfrak{t}'\)（\(\mathfrak{m}'\) 的中心）由半单元素组成。则 \(\mathfrak{r} = \mathfrak{t} \oplus \mathfrak{n}_{\mathrm{V}}(\mathfrak{g}) = \mathfrak{t}' \oplus \mathfrak{n}_{\mathrm{V}}(\mathfrak{g})\)。鉴于命题 6 的推论 2，我们化为情形 \(\mathfrak{t} = \mathfrak{t}'\)。则 \(\mathfrak{s}' \subset \mathfrak{z}\)；由于 \(\dim \mathfrak{s}' = \dim \mathfrak{s}\)，\(\mathfrak{s}'\) 是 \(\mathfrak{z}\) 的一个 Levi 子代数。由第一章 §6 第 8 节定理 5，存在 \(x \in \mathfrak{n}_{\mathrm{V}}(\mathfrak{z})\)，使得 \(e^{\operatorname{ad} x}(\mathfrak{s}) = \mathfrak{s}'\)；由于 \(x\) 与 \(\mathfrak{t}\) 交换，还有 \(e^{\operatorname{ad} x}(\mathfrak{t}) = \mathfrak{t}\)。

